\subsection{Experimental Setup}

\paragraph{Dataset.}
We use the CMCC corpus, which contains texts from 10 authors across 5 genres (blogs, chat, discussion, emails, essays) and 6 topics, yielding up to 30 (genre, topic) writing tasks per author.

\paragraph{Conditions.}
We evaluate generalization under three increasingly challenging data splits, each constructed independently per author:
\begin{itemize}
    \item \textbf{Seen}: A random 8-example train / rest-test split (3 seeds).
    \item \textbf{Unseen-Genre}: All examples from one held-out genre are reserved for testing; the model trains on the remaining 4 genres. We repeat this for each of the 5 genres.
    \item \textbf{Unseen-Topic}: All examples from a pair of held-out topics go to test (3 topic-pair variants: \{Iraq, Gender Discrimination\}, \{Catholic Church, Gay Marriage\}, \{Privacy, Marijuana\}).
\end{itemize}

\paragraph{Training.}
For each condition $\times$ author, we train a separate \textsc{Ditto} model on Mistral-7B-Instruct-v0.2 using LoRA ($r=32$, $\alpha=64$). Training proceeds in two phases: SFT on author demonstrations ($\leq$30 steps, early-stopped at loss $< 1.0$), followed by iterative DPO with online comparison generation (40 steps, mixing 70\% expert-vs-current, 20\% expert-vs-replay, 10\% inter-policy pairs).

\paragraph{Generation \& Evaluation.}
We generate 3 samples per test prompt from each trained model. A GPT-4o judge rates each generation against the author's ground-truth reference on a 1--7 Likert scale for stylistic similarity, where 1 = extremely different and 7 = extremely similar.
